\subsubsection{LLMPage} 
\label{sec:LLMPage}
\paragraph{Descrizione}
LLMPage è il componente che si occupa di rendere visibile le opzioni e informazioni necessarie creazione e modifica di un LLM  con le seguenti funzionalità:
\begin{itemize}
    \item comportamento modulabile in fase di routing a seconda della chiamata HTTP, consentendo comportamenti specifici a seconda della chiamata; ciò consente ad un componente di poter gestire leggere variazioni di comportamento conservando la stessa logica di fondo. Tale parametro consente la differenziazione tra la funzione di visualizzare gli LLM e quella di visualizzare gli LLM in fase di test.
    \item Modifica del nome dell'LLM.
    \item Modifica del URL dell'LLM.
    \item Modifica della chiave associata alla domanda dell'LLM.
    \item Modifica della chiave associata alla risposta dell'LLM.
    \item Modifica di una coppia chiave-valore associata alla creazione dell'header dell'LLM.
    \item Modifica di una coppia chiave-valore associata alla creazione del body dell'LLM.
    \item Salvataggio del nome dell'LLM.
    \item Salvataggio del URL dell'LLM.
    \item Salvataggio della chiave associata alla domanda dell'LLM.
    \item Salvataggio della chiave associata alla risposta dell'LLM.
    \item Salvataggio di una coppia chiave-valore associata alla creazione dell'header dell'LLM.
    \item Salvataggio di una coppia chiave-valore associata alla creazione del body dell'LLM.
    \item Gestione per la conferma di azioni che comporteranno modifiche significative.
\end{itemize}   
\paragraph{SottoComponenti}
LLMPage è composto dai seguenti sotto-componenti:
\begin{itemize}        
        \item \textbf{LLMFormComponent}: Componente che si occupa del form per la creazione e modifica degli LLM.
        \item \textbf{FieldReqComponent}: Tale componente funge da wrapper per la chiamata HTTP di salvataggio per la conversione del tipo ritornato dalla chiamata, in un tipo consono alla logica del componente LLMPage.
    \end{itemize}

\paragraph{Tracciamento dei requisiti per LLMPage}
\begin{tabularx}{\textwidth}{|c|c|}
    \hline
    \textbf{ID} & \textbf{componente} \\
    \hline
        RFF-4.01 & LLMFormComponent \\
    \hline
        RFF-4.02 & LLMFormComponent \\
    \hline
        RFF-4.03 & LLMFormComponent \\
    \hline
        RFF-4.04 & LLMFormComponent \\
    \hline
        RFF-4.05 & LLMFormComponent \\
    \hline
        RFF-4.06 & LLMFormComponent \\
    \hline
        RFF-4.07 & LLMFormComponent \\
    \hline
        RFF-4.08 & LLMFormComponent \\
    \hline
        RFF-4.09 & FieldReqComponent \\
    \hline
        RFF-4.10 & LLMFormComponent \\
    \hline
        RFF-4.11 & MessageComponent \\
    \hline
        RFF-4.12 & MessageComponent \\
    \hline
        RFF-4.13 & MessageComponent \\
    \hline
        RFF-5.01 & LLMFormComponent \\
    \hline
        RFF-5.02 & LLMFormComponent \\
    \hline
        RFF-5.03 & LLMFormComponent \\
    \hline
        RFF-5.04 & LLMFormComponent \\
    \hline
        RFF-5.05 & LLMFormComponent \\
    \hline
        RFF-5.06 & LLMFormComponent \\
    \hline
        RFF-5.07 & FieldReqComponent \\
    \hline
        RFF-5.08 & LLMFormComponent \\
    \hline
        RFF-5.09 & MessageComponent \\
    \hline
        RFF-5.10 & MessageComponent \\
    \hline
        RFF-5.11 & MessageComponent \\
    \hline
\end{tabularx}